Пусть $\displaystyle ( \Omega ,\mathcal{F})$ и $\displaystyle ( E,\mathcal{E})$ -- измеримые пространства.
\begin{definition}
	Если $\displaystyle \xi :\Omega \rightarrow E$ таково, что $\displaystyle \xi ^{-1}( B) \in \mathcal{F} \ \forall B\in \mathcal{E}$. Тогда $\displaystyle \xi $ называется \textit{случайным элементом}. Если $\displaystyle ( E,\mathcal{E}) =\left(\mathbb{R}^{m} ,\mathcal{B}\left(\mathbb{R}^{m}\right)\right)$, то $\displaystyle \xi $ -- \textit{случайный вектор}. 
\end{definition}
Пусть $\displaystyle ( \Omega ,\mathcal{F} ,P)$ -- вероятностное пространство.
\begin{definition}
	\textit{Распределением} случайного элемента $\displaystyle \xi $ называют меру $\displaystyle P_{\xi }$ на $\displaystyle \mathcal{E}$, такую что $\displaystyle P_{\xi }( B) =P( \xi \in B)$.
\end{definition}

Пусть наблюдается некоторый эксперимент над $\displaystyle m$-мерным случайным вектором с распределением $\displaystyle P$. Нужно построить математическую модель эксперимента, то есть построить вероятностное пространство и случайный вектор над этим вероятностым пространством с распределением $\displaystyle P$.
\begin{definition}
	Множество $\displaystyle \mathcal{X}$ всевозможных исходов эксперимента называется \textit{выборочным пространством.}
\end{definition}
    Определим $\displaystyle \mathcal{B}_{\mathcal{X}}$ -- некоторую $\displaystyle \sigma $-алгебру над выборочным множеством (обычно будем считать ее борелевской) и вероятностную меру $\displaystyle P$ на измеримом пространстве $\displaystyle (\mathcal{X} ,\mathcal{B}_{\mathcal{X}})$.
    
    Также, определим функцию $\displaystyle X( x) =x,\ \forall x\in \mathcal{X}$.
\begin{proposition}
	Функция $\displaystyle X:\mathcal{X}\rightarrow \mathcal{X}$ является \textit{случайным элементом} на вероятностном пространстве $\displaystyle (\mathcal{X} ,\mathcal{B}_{\mathcal{X}} ,P)$ и имеет распределение $\displaystyle P_{X} =P$.
\end{proposition}
\begin{proof}
    \begin{gather*}
        \forall B\in \mathcal{B}_{\mathcal{X}} \ X^{-1}( B) =\{x\in \mathcal{X} :X( x) \in B\} =\{x\in \mathcal{X} :x\in B\} = B.
    \end{gather*}
    Значит, $X$ -- случайный элемент, и
    \begin{gather*}
        \forall B\in \mathcal{B}_{\mathcal{X}} \ P_{X}( B) =P( X\in B) =P( x\in \mathcal{X} :X( x) \in B) =P( x\in \mathcal{X} :x\in B) =P( B).
    \end{gather*}
\end{proof}

\begin{definition}
	Фукнция $\displaystyle X$, определенная выше, называется \textit{наблюдением}.
\end{definition}
Построим модель $n$ независимых экспериментов. Пусть 
\begin{gather*}
    \mathcal{X}^{n} =\mathcal{X} \times \dotsc \times \mathcal{X},\\ \mathcal{B}_{\mathcal{X}}^{n} =\mathcal{B}\left(\mathcal{X}^{n}\right) =\sigma ( B_{1} \times \dotsc \times B_{n}, B_{i} \in \mathcal{B}_{\mathcal{X}}),\\ P^{n}( B_{1} \times \dotsc \times B_{n}) =P( B_{1}) \cdotp \dotsc \cdotp P( B_{n}).
\end{gather*}
Снова рассмотрим тождественное отображение $\displaystyle X:\mathcal{X}^{n}\rightarrow \mathcal{X}^{n}$. Также, определим функцию $\displaystyle X_{i} :\mathcal{X}^{n}\rightarrow \mathcal{X},\ X_{i}( x_{1} ,\dotsc ,x_{n}) =x_{i}$.
\begin{proposition}
	$\displaystyle X_{i}$ -- случайная величина на вероятностном пространстве $\displaystyle \left(\mathcal{X}^{n} ,\mathcal{B}_{\mathcal{X}}^{n} ,P^{n}\right)$ с распределением $\displaystyle P$, а $\displaystyle X_{1} ,\dotsc ,X_{n}$ -- независимы в совокупности.
\end{proposition}
\begin{proof}
    \begin{gather*}
        \forall B\in \mathcal{B}_{\mathcal{X}} \ X^{-1}( B) =\left\{\overline{x} \in \mathcal{X}^{n} :X(\overline{x}) \in B\right\} =\left\{\overline{x} \in \mathcal{X}^{n} :x_{i} \in B\right\} =\\ \mathcal{X} \times \dotsc \times \mathcal{X} \times B\times \mathcal{X} \times \dotsc \times \mathcal{X} \in \mathcal{X}^{n}.
    \end{gather*}

    Зафиксируем 
    $\displaystyle i\in \{1,\dotsc ,n\}$. Тогда
    \begin{gather*}
        P^{n}( X_{i} \in B_{i}) =P^{n}\left(\left\{\overline{x} \in \mathcal{X}^{n} :X_{i}(\overline{x}) \in B_{i}\right\}\right) =P^{n}\left(\left\{\overline{x} \in \mathcal{X}^{n} :x_{i} \in B_{i}\right\}\right) =\\ P^n(\mathcal{X} \times \dotsc \times \mathcal{X} \times B_{i} \times \mathcal{X} \times \dotsc \mathcal{X}) =\ P(\mathcal{X}) \cdot \dotsc \cdot P(\mathcal{X}) \cdot P( B_{i}) \cdot P(\mathcal{X}) \cdot \dotsc \cdot P(\mathcal{X}) =P( B_{i}).
    \end{gather*}
    Докажем, что случайные величины независимы в совокупности.
    \begin{gather*}
        P^{n}( X_{1} \in B_{1} ,\dotsc ,\ X_{n} \in B_{n}) =P^{n}\left(\overline{x} \in \mathcal{X}^{n} :X_{1}(\overline{x}) \in B_{1} ,\dotsc ,\ X_{n}(\overline{x}) \in B_{n}\right) =\\ P^n\left(\overline{x} \in \mathcal{X}^{n} :x_{1} \in B_{1} ,\dotsc ,\ x_{n} \in B_{n}\right) =P^{n}( B_{1} \times \dotsc \times B_{n}) =P( B_{1}) \dotsc P( B_{n}) =\\ P^n( X_{1} \in B_{1}) \dotsc P^n( X_{n} \in B_{n}).
    \end{gather*}
\end{proof}

\begin{definition}
	Совокупность $\displaystyle X=( X_{1} ,\dotsc ,X_{n})$ независимых одинаково распределенных случайных векторов с распределением $\displaystyle P$ называется выборкой размера $\displaystyle n$ с распределением $\displaystyle P$.
\end{definition}
Часто нам придется рассматривать выборку при $\displaystyle n\rightarrow \infty $. Поэтому, введем следующие объекты: 
\begin{gather*}
    \mathcal{X}^{\infty } =\mathcal{X} \times \mathcal{X} \times \dotsc =\{( x_{1} ,x_{2} ,\dotsc ) :x_{i} \in \mathcal{X}\},\\ \displaystyle \mathcal{B}_{\mathcal{X}}^{\infty } =\sigma \left(\left\{( x_{1} ,x_{2} ,\dotsc ) :\exists n\in \mathbb{N} ,\ \exists B\in \mathcal{B}_{\mathcal{X}}^{n} \hookrightarrow ( x_{1} ,\dotsc ,x_{n}) \in B\right\}\right),
\end{gather*}
$\displaystyle P^{\infty }$ -- вероятностная мера на $\displaystyle \left(\mathcal{X}^{\infty } ,\mathcal{B}_{\mathcal{X}}^{\infty }\right)$, такая что 
\begin{gather*}
    P^{\infty }( B_{1} \times B_{2} \times \dotsc \times B_{n} \times \mathcal{X} \times \mathcal{X} \times \dotsc ) =P( B_{1}) \dotsc P( B_{n}).
\end{gather*}
Снова определим $\displaystyle X:\mathcal{X}^{\infty }\rightarrow \mathcal{X}^{\infty }$ -- тождественное отображение, $\displaystyle X_{i}( x_{1} ,x_{2} ,\dotsc ) =x_{i}$. Тогда $\displaystyle X_{1},\ X_{2} ,\dotsc $ -- независимы и одинаково распределены с распределением $\displaystyle P$.
\\
В дальнейшем для простоты обозначений будем писать $\displaystyle (\mathcal{X},\ \mathcal{B}_{\mathcal{X}},\ P)$ вместо $$\displaystyle \left(\mathcal{X}^{n},\ \mathcal{B}_{\mathcal{X}}^{n},\ P^{n}\right)$$ и \ $$\displaystyle \left(\mathcal{X}^{\infty },\ \mathcal{B}_{\mathcal{X}}^{\infty },\ P^{\infty }\right)$$

Пусть $\displaystyle \mathcal{P}$ -- семейство вероятностных мер на измеримом пространстве $\displaystyle (\mathcal{X} ,\mathcal{B}_{\mathcal{X}})$.
\begin{definition}
	Тройка $\displaystyle (\mathcal{X},\ \mathcal{B}_{\mathcal{X}},\ \mathcal{P})$ называется \textit{вероятностно-статистической моделью}.
\end{definition}
\begin{definition}
	$\displaystyle X_{1}( \omega ) =x_{1},\,\ldots,\, X_{n}(\omega) = x_{n}$ будем называть\textit{ реализацией выборки}.
\end{definition}
\begin{note}
Задача статистики - сделать вывод о неизвестном распределении по реализции выборки 
\end{note}
\textbf{Параметрическая модель}
\begin{definition}
Вероятностно-статистическая модель называется параметрической, если семейство распределений $\displaystyle \mathcal{P}$ параметризовано, т.е. $\displaystyle \mathcal{P} =\{P_{\theta } :\theta \in \Theta \}$.
\end{definition}
\begin{example}
$\displaystyle \mathcal{P} =\{exp( \theta ),\ \theta  >0\}$ -- семейство экспоненциальных распределений, параметризованных $\displaystyle \theta $.
\end{example}